% ===================================================================
%  TAULA N° 15 — Lo Mercat. — Los Aliments.
%  Traduction 2026 d'apres texto/io, controlee sur texto/fr et
%  texto/ca. Consigne : voir texto/oc/10-taula-01.tex.
% ===================================================================

%%K t15-apar-1 apar
\VUsaut{4.0mm}
\VUtitre{15.0pt}{60}{TAULA N° 15}
\VUsaut{3.0mm}

%%K t15-tit-1 sub
\VUpk{11.4pt}{40}{Lo Mercat. --- Los Aliments.}
\VUsaut{3.0mm}

%%K t15-01-1 p
1. --- La màger part dels comerçants que nos provesisson dels aliments
necessaris son acampats jos la \VUgras{nau} \textsuperscript{(1)} del \VUgras{mercat
cobèrt}, o dins los carrièrs vesins.

%%K t15-02-1 p
2. --- Lo \VUgras{cansaladièr} \textsuperscript{(2)}, que vend de \VUgras{cambajon}
\textsuperscript{(3)}, de \VUgras{bodin} \textsuperscript{(4)}, de
\VUgras{salsissòt} \textsuperscript{(5)}, de \VUgras{gòfa} \textsuperscript{(6)} e
de \VUgras{cansalada} \textsuperscript{(7)}, es contra la \VUgras{vendeirèla de
burre e de fromatge} \textsuperscript{(8)}, que son estal es cobèrt de
\VUgras{fromatges de bòla} \textsuperscript{(9)} amb una crosta roja, de
\VUgras{camemberts} \textsuperscript{(10)}, de grandas \VUgras{ròdas de gruèire}
\textsuperscript{(11)} e de \VUgras{pans de burre} \textsuperscript{(12)} fresc.

%%K t15-03-1 p
3. --- Davant ela, una \VUgras{legumièira} \textsuperscript{(13)} ofrís a sos
clients de bèles \VUgras{tomatis} \textsuperscript{(14)}, de \VUgras{caulsflors}
\textsuperscript{(15)}, d'\VUgras{ensaladas} \textsuperscript{(16)}, de
\VUgras{cauls} \textsuperscript{(17)}, de \VUgras{cocorletas}
\textsuperscript{(19)}, de \VUgras{melons} \textsuperscript{(18)} e quitament de
\VUgras{bolets} \textsuperscript{(20)}. Mas un \VUgras{cerveseire}, qu'a arrestat
just davant son estal son \VUgras{carri de liurason} \textsuperscript{(21)} cargat
de \VUgras{barricas de cervesa} \textsuperscript{(22)}, empacha sos clients de
s'apropchar. Un ortolan li pòrta un panièr de \VUgras{carchòfas}
\textsuperscript{(23)}.

%%K t15-tit-2 sub
\VUsaut{4.0mm}
\VUpk{10.2pt}{40}{Vendre e crompar.}
\VUsaut{3.0mm}

%%K t15-04-1 p
4. --- Al mercat l'animacion es granda, qu'es arribada l'ora de l'\VUgras{encant}
\textsuperscript{(24)}. Las \VUgras{mestressas de l'ostal} \textsuperscript{(26)}
que venon aicí far la crompa s'arrèstan en passar davant l'estal de la
\VUgras{tripièira} \textsuperscript{(25)} per crompar de \VUgras{tripas} o un
\VUgras{cap de vedèl}.

%%K t15-05-1 p
5. --- Un \VUgras{cosinièr} \textsuperscript{(27)} gròs marchanda una
\VUgras{tonina} fòrça bèla amb la \VUgras{peissonièira} \textsuperscript{(28)},
qu'auça los prètzes a son grat. A recebut una granda partida d'\VUgras{anguilas,
solas, truchas} e de \VUgras{carpas}. Son \VUgras{merluç} e sos \VUgras{muscles}
son fòrça fresques.

%%K t15-tit-3 sub
\VUsaut{4.0mm}
\VUpk{10.2pt}{40}{Lo barat e lo car.}
\VUsaut{3.0mm}

%%K t15-06-1 p
6. --- La \VUgras{polalhièira} \textsuperscript{(29)}, installada a son costat, es
fòrça onèsta. Quand vend un \VUgras{piòt}, una \VUgras{auca}, un \VUgras{anet} o
un simple \VUgras{polet}, demanda pas mai que lo prètz just.

%%K t15-07-1 p
7. --- Al canton de la plaça, al primièr estatge, una \VUgras{linjaira}
\textsuperscript{(30)} s'i es establida i a pas gaire, ont las crompairas son
totjorn seguras de trobar un assortiment fòrça bon de \VUgras{toalhas},
\VUgras{serrietas, davantals} e \VUgras{draps de cosina}. Al meteis estatge, un
\VUgras{cotelièr} \textsuperscript{(31)} a montat son obrador. Agusa los
\VUgras{cotèls} de las legumièiras e fa als ortolans de \VUgras{podadoiras} de
l'\VUgras{acièr} mai dur.

%%K t15-tit-4 sub
\VUsaut{4.0mm}
\VUpk{10.2pt}{40}{La botiga de viures.}
\VUsaut{3.0mm}

%%K t15-08-1 p
8. --- Jos son obrador s'es dobèrta i a pas gaire una granda botiga de viures. Es
ja pas mai l'\VUgras{espiçariá} \textsuperscript{(32)} simpla e modèsta dels temps
passats, que vendiá pas que de gèneres corrents --- \VUgras{tè}
\textsuperscript{(33)}, \VUgras{chocolat} \textsuperscript{(34)}, \VUgras{pan de
sucre} \textsuperscript{(37)} e \VUgras{sabon} \textsuperscript{(38)}. Aicí se
vend de tot: de \VUgras{gaufretas crocantas}, de \VUgras{consèrvas}
\textsuperscript{(35)}, de \VUgras{licors} \textsuperscript{(36)}, de
\VUgras{rom, conhac, kirsch, anís} e de \VUgras{curaçao}. S'i tròba de caça amb la
meteissa aisidesa --- \VUgras{lèbre} \textsuperscript{(39)}, \VUgras{tordres}
\textsuperscript{(40)}, \VUgras{perdises} \textsuperscript{(41)},
\VUgras{calhas} \textsuperscript{(42)}, \VUgras{anets sauvatges}
\textsuperscript{(43)} o \VUgras{faisan} \textsuperscript{(44)} --- coma de
frucha tropicala tala coma de \VUgras{banans} \textsuperscript{(46)} e
d'\VUgras{ananàs}.

%%K t15-09-1 p
9. --- Un \VUgras{comés} \textsuperscript{(45)} fòrça servicial es prèst a liurar
çò que se li demande: de \VUgras{confitura} \textsuperscript{(47)}, de grasèlas o
de codonhs, de \VUgras{saïm} \textsuperscript{(48)}, de \VUgras{cogombrets}
\textsuperscript{(49)} o de \VUgras{sardinas} \textsuperscript{(50)} en saladura.

%%K t15-09-2 p
Aquò es fòrça comòde per las \VUgras{practicas} \textsuperscript{(51)}, que
tròban, al costat dels gèneres mai corrents, las primícias mai raras e la frucha
mai saborosa: de \VUgras{peras} \textsuperscript{(52)} sucosas, de
\VUgras{prunas} \textsuperscript{(53)}, de \VUgras{rasims}
\textsuperscript{(54)}, de \VUgras{persègas} \textsuperscript{(55)},
d'\VUgras{amètlas} \textsuperscript{(56)} e de \VUgras{noses}
\textsuperscript{(57)}.

%%K t15-tit-5 sub
\VUsaut{4.0mm}
\VUpk{10.2pt}{40}{Las practicas.}
\VUsaut{3.0mm}

%%K t15-10-1 p
10. --- Lo \VUgras{ris} \textsuperscript{(58)} e las \VUgras{lentilhas}
\textsuperscript{(59)} son contra la caissa d'\VUgras{arencs}
\textsuperscript{(60)} fumats; la \VUgras{trufa} \textsuperscript{(61)} e la
\VUgras{monjeta} \textsuperscript{(63)} se tòcan amb los \VUgras{poms}
\textsuperscript{(62)}, las \VUgras{figas} \textsuperscript{(64)}, las
\VUgras{prunas secas} \textsuperscript{(65)} e las \VUgras{avelanas}
\textsuperscript{(66)}. Se tròban tanben dins aquesta botiga d'\VUgras{uòus}
\textsuperscript{(67)} fresques e d'\VUgras{olivas} \textsuperscript{(68)}, de
biais que los \VUgras{panièrs} \textsuperscript{(70)} e los \VUgras{filats de la
crompa} \textsuperscript{(73)} s'emplenan a tota pressa.

%%K t15-11-1 p
11. --- Lo \VUgras{cafè} \textsuperscript{(72)} d'aquesta excellenta maison a
ganhat una fama plan meritada entre las \VUgras{servicialas}
\textsuperscript{(69)} e las cosinièiras, que lo vièlh \VUgras{espicièr}
\textsuperscript{(71)} lo torrifica el meteis amb lo mai grand suènh.

%%K t15-tit-6 sub
\VUsaut{4.0mm}
\VUpk{10.2pt}{40}{Causas de veire.}
\VUsaut{3.0mm}

%%K t15-12-1 p
12. --- Tot lo mond ven pas al mercat crompar de provisions. Dins lo pintoresc anar
e venir del carrieròl que mena a la nau, se vei lo mond mai variat. Contra un
brave \VUgras{tuaire} \textsuperscript{(75)} que buta davant el una \VUgras{bèrla}
\textsuperscript{(74)}, vaicí dos soldats en permission, fòrça fièrs de mostrar
lors unifòrmes vistoses: l'\VUgras{ussard} \textsuperscript{(76)} amb son
\VUgras{dolman} blau e son \VUgras{chacò a plumet}, e lo \VUgras{caporal de zoaus}
amb sas bragas largas e un \VUgras{fez} sul cap.

%%K t15-13-1 p
13. --- Lo \VUgras{fornièr} \textsuperscript{(80)}, drech sul lindal del
\VUgras{forn} \textsuperscript{(78)}, ven de pastar la pasta de son \VUgras{pan}
\textsuperscript{(79)} e de traire del forn los \VUgras{croissants}
\textsuperscript{(81)}, los \VUgras{panets} \textsuperscript{(82)} e las
\VUgras{fogaças} \textsuperscript{(83)} que son dins la vitrina.

%%K t15-14-1 p
14. --- Dessús lo forn demòra una abila \VUgras{cordurièira}
\textsuperscript{(84)} qu'emplega mantuna \VUgras{brodaira}
\textsuperscript{(85)}. A son costat demòra una \VUgras{bugadièira e
planchaira} \textsuperscript{(86)}, qu'empesa la \VUgras{linja}
\textsuperscript{(88)} aprèp l'aver lavada e eissugada, e la plancha amb una
\VUgras{plancha} \textsuperscript{(87)}.

%%K t15-tit-7 sub
\VUsaut{4.0mm}
\VUpk{10.2pt}{40}{Carn, peis e legumas.}
\VUsaut{3.0mm}

%%K t15-15-1 p
15. --- Dins la \VUgras{maselariá} \textsuperscript{(89)} del ras-de-caussada se
vend de carn de tota mena --- \VUgras{moton} \textsuperscript{(90)},
\VUgras{buòu} \textsuperscript{(94)} o \VUgras{vedèl} \textsuperscript{(92)} ---
en \VUgras{cuèissas d'anhèl, bistècas, rostits} e \VUgras{costeletas}. Lo
\VUgras{varlet de masèl} \textsuperscript{(93)} talha tota aquela carn e met de
costat los \VUgras{òsses de mesolha} \textsuperscript{(96)}. A la pòrta de la
maselariá i a una \VUgras{ostrièira} \textsuperscript{(97)} que vend d'\VUgras{ostras}
\textsuperscript{(98)}. Las dobrís se se li demanda.

%%K t15-16-1 p
16. --- Totes aqueles comerçants pagan una patenta o un drech d'estal; per aquò lo
\VUgras{garda} \textsuperscript{(100)} perseguís sens pietat las pauras
\VUgras{legumièiras ambulantas} \textsuperscript{(99)}, que lor fan la
concurréncia en portar dins lors bèrlas de \VUgras{caròtas, rabes, naps, pòrres}
o de \VUgras{manats d'espargues, pesets fresques, espinarcs, agrelas, creissons} e
de \VUgras{jorverd}.

%%K t15-tit-8 sub
\VUsaut{4.0mm}
\VUpk{10.2pt}{40}{Mèfi al garda!}
\VUsaut{3.0mm}

%%K t15-17-1 p
17. --- Lo dròlle que vend d'\VUgras{alhs} \textsuperscript{(101)} e de
\VUgras{cebas} sap fòrça plan qu'a el tanpauc li es pas permés de vendre sa
mercadariá contra lo mercat.

%%K t15-17-1 p suite
Se met a córrer, amb lo risc de tombar lo panièr de \VUgras{crancs}, de
\VUgras{escamarlans} e de \VUgras{gambetas} de la \VUgras{vendeirèla}
\textsuperscript{(102)} de \VUgras{langostas} e de \VUgras{lingostins}.

%%K t15-18-1 p
18. --- Tot lo mond s'afana e fa un grand bruch; mas aquò empacha pas d'ausir
clarament lo crit del \VUgras{veirièr} \textsuperscript{(103)} ambulant, ni lo
crit del \VUgras{carbonièr} \textsuperscript{(104)}, que ven de daissar un moment
son carri per liurar un \VUgras{sac de carbon} \textsuperscript{(105)}.
\VUsaut{6.0mm}
\VUfilet{22mm}
